\chapter{Problem Statement}
\label{ch:problem}

\section{Problem Definition}

Given one panoramic radiograph $I$ of a patient, the system must produce a report $R(I)$ with one
record per tooth present in the image. Each record $t_i$ carries:
\begin{itemize}
\item the tooth's box in image coordinates and the detector's confidence;
\item its FDI code $c_i \in \{11,\dots,18,\,21,\dots,28,\,31,\dots,38,\,41,\dots,48\}$ and the
  numbering confidence;
\item a caries finding in $\{\textit{no caries},\ \textit{caries},\ \textit{deep caries}\}$ with its
  confidence;
\item an abnormality flag (the Stage-4 output) with its confidence;
\item an \emph{assessable} flag that says whether the record is reliable enough to read.
\end{itemize}
The report must be produced on a CPU in a few seconds, must state its own limitations, and must be
traceable to the exact model checkpoints that produced it. The problem is therefore not a single
classification but a structured prediction: which teeth are present, which tooth each one is, and
what each one shows.

\section{The Report as a Chain of Decisions}

A record is correct only if the tooth is detected ($D$), numbered correctly once detected ($N$), and
given the right finding ($F$):
\begin{equation}
P(\text{record correct}) = P(D)\;P(N \mid D)\;P(F \mid D, N).
\label{eq:chain}
\end{equation}
The first two factors concern identity, not disease, and bound everything after them
(Figure~\ref{fig:errorchain}). A perfect caries model still writes a wrong report on every tooth that
was missed or misnumbered: a lesion correctly recognised but attached to the wrong tooth number is a
wrong report. A score that summarises only the last decision is silent about the two that condition
it. On the DENTEX evaluation images the first two factors are 99.6\% and 93.2\%, so 92.8\% of teeth
reach the third decision correctly identified (Table~\ref{tab:identity}). % {C030}

\begin{figure}[H]
\centering
% Error chain of a tooth-indexed report (eq. 2.1). Values traced to paper/paper.md Table 1
% (DENTEX 164 images, 936 teeth: detected 99.6%, exact of detected 93.2%, end to end 92.8%).
\begin{tikzpicture}[node distance=5mm and 9mm]
\node[io,text width=24mm] (in) {OPG\\[-1pt]{\normalfont one radiograph}};
\node[stage,text width=30mm,right=of in] (d) {Detected?\\[-1pt]{\normalfont $P(D)$ = 99.6\%}};
\node[stage,text width=30mm,right=of d] (n) {Numbered correctly?\\[-1pt]{\normalfont $P(N\mid D)$ = 93.2\%}};
\node[cari,text width=30mm,right=of n] (f) {Right finding?\\[-1pt]{\normalfont $P(F\mid D,N)$}};
\node[rept,text width=24mm,below=13mm of f] (ok) {Correct record};
\draw[arr] (in) -- (d); \draw[arr] (d) -- node[lbl,above]{yes} (n); \draw[arr] (n) -- node[lbl,above]{yes} (f);
\draw[arr] (f) -- node[lbl,right]{yes} (ok);
\node[soft=p2maroon,text width=28mm,below=13mm of d] (e1) {Tooth absent\\from the report};
\node[soft=p2maroon,text width=28mm,below=13mm of n] (e2) {Findings attached\\to the wrong tooth};
\draw[arr,draw=p2maroon] (d) -- node[lbl,left]{no} (e1);
\draw[arr,draw=p2maroon] (n) -- node[lbl,left]{no} (e2);
\node[lbl,above=1mm of n,xshift=-20mm] {identity: 92.8\% of teeth reach the finding correctly identified};
\end{tikzpicture}

\caption{The error chain of a tooth-indexed report. An error in an early decision cannot be repaired
by a later one. Percentages: the pipeline on 164 held-out DENTEX images (936 teeth), from
Table~\ref{tab:identity}.}
\label{fig:errorchain}
\end{figure}

\section{Why the Usual Caries Score Can Mislead}

Caries is not spread evenly over the mouth. In the project's data the first and second molars carry
most of it (Figure~\ref{fig:prior}). A learner trained to minimise error over a population will use
any cue that predicts the label, and tooth position predicts it, so a model can earn part of its
score by knowing \emph{where caries usually is} rather than seeing where it is in this image, a
shortcut in the sense of Geirhos et al.~\cite{geirhos2020shortcut}. The usual metric cannot tell the
two apart, because a set of tooth codes can be chosen without any image.

Two consequences follow. First, a reported score has a floor set by this prior, and it means little
until an image-blind constant is scored beside it. Second, the floor moves with the population: a
constant's lesion agreement is the share of a cohort's lesions that lie on its teeth, so agreement
measured on one development cohort cannot separate a model that reads the lesion from one that has
absorbed the prior. External cohorts, paired with the constant and reported per site, can.
Chapter~\ref{ch:method} derives the best image-blind constant in closed form.

\section{Research Questions}

\begin{description}[leftmargin=1.2cm,style=nextline]
\item[RQ1] Does tooth-level agreement separate the full pipeline from an image-blind constant, on
  development data and on external cohorts?
\item[RQ2] Does changing the caries stage from classifying tooth crops to localising lesions in the
  whole image, a task a fixed list cannot perform, separate them?
\item[RQ3] What does the pipeline show on outputs a fixed list cannot produce: numbering, calibrated
  confidence and abstention?
\end{description}

\section{Constraints and Challenges}

\begin{itemize}
\item \textbf{No single dataset carries every label.} Tooth boxes and numbers, caries labels,
  lesion boxes and abnormality masks come from different public datasets with different conventions,
  so each stage is trained on the data that carries its label (Table~\ref{tab:datasets}).
\item \textbf{Labels are not diagnoses.} DENTEX caries labels were read from panoramic radiographs
  alone, and the Tufts abnormality masks mark apical lesions, sclerosis and cysts, not periodontal
  bone loss.
\item \textbf{Evaluation sets are small and clustered.} Teeth within one radiograph share exposure
  and annotator, so intervals must resample radiographs, not teeth.
\item \textbf{Domain shift.} Radiographs from another hospital differ in device, exposure and
  population; one external cohort (Ningxia) is kept locked for testing.
\item \textbf{Deployment.} Inference must run on a CPU, and the report must carry its own caveats.
\item \textbf{Licences.} DENTEX images may be shown under CC~BY-NC-SA~4.0; Tufts images may not be
  redistributed.
\end{itemize}

\section{Success Criteria}

A change to the pipeline is accepted only under a rule written before the evidence. For the caries
stage the rule, committed before any candidate was trained, required the new method to beat both
the image-blind constant and the existing crop classifier on at least two of three cohorts, with the
interval excluding zero, at no more than 1.5 times their rate of flagged teeth
(Section~\ref{sec:method-prereg}). % {C020} {C021}
Every reported number must trace to a saved artefact, and a
result that fails is reported beside one that passes.
